\section{The \ose{} Open-Source Ecosystem}

% REQUIRED by the solicitation for EVERY PESOSE proposal:
%   "details on the current status of the open-source product development and testing model,
%    methods of dissemination, user base, and contributor base."
% This section supplies exactly that for FreeRTOS. Keep all URLs out -> cite \cite{freertos} etc.

\noindent\textbf{The product.}
% TODO(pesose): what FreeRTOS is -- a real-time kernel + libraries (FreeRTOS+TCP, FreeRTOS+FAT,
% coreMQTT, corePKCS11, etc.); MIT license; portable across many MCU architectures.
\ose{} is an MIT-licensed real-time operating system kernel for microcontrollers and small microprocessors, together with a set of companion libraries~\cite{freertos}.

\noindent\textbf{Development and testing model.}
% TODO(pesose): how FreeRTOS is developed and tested today -- GitHub-based, AWS stewardship,
% maintainers, release cadence, existing CI, existing test/coverage practices, MISRA conformance,
% existing static-analysis usage. Be accurate; cite sources in References Cited.

\noindent\textbf{Methods of dissemination.}
% TODO(pesose): GitHub repositories, official site, package ecosystems, vendor SDKs that bundle
% FreeRTOS (e.g., chip-vendor toolchains), Amazon-FreeRTOS / FreeRTOS LTS distributions.

\noindent\textbf{User base.}
% TODO(pesose): scale and breadth of deployment (billions of devices; sectors: industrial control,
% medical, automotive, consumer IoT, aerospace). Quantify with citable figures where possible.

\noindent\textbf{Contributor base.}
% TODO(pesose): who develops FreeRTOS -- AWS team + external contributors, GitHub contributor counts,
% governance (note: this informs the upstream-adoption strategy in Activity 3).

\noindent\textbf{The problem being addressed.}
% TODO(pesose): crisp statement of the safety/security/privacy problem in the FreeRTOS ecosystem
% that this effort will reduce (leads into the Vulnerabilities & Risks section).
